%!TEX root = main.tex

\begin{abstract}

Broadcast is a fundamental operation in Mobile Ad-Hoc Networks (MANETs).
A large variety of broadcast algorithms have been proposed. 
They differ in the way message forwarding between nodes is controlled\deleted{ in time and space}, and in the level of information about the topology that this control requires.
%
Deployment scenarios for MANETs vary widely, in particular in terms of nodes density and mobility.
%
The choice of an\deleted{broadcast} algorithm depends on its expected reachability and energy cost, which are both\deleted{ highly} impacted by the deployment context.
%
In this work, we are interested in overcoming the lack of comprehensive comparison of the costs and effectiveness of broadcast algorithms for MANETs, and ease an informed choice between them based on target environmental conditions.
%
We describe the results of an\deleted{ independent} experimental study of five algorithms, representative of the main design alternatives.
%
Our study reveals that the best algorithm for a given situation, such as a high density and a stable network, is not necessarily the most appropriate for a different situation such as a sparse and mobile network.
We identify the algorithms characteristics that are correlated with these differences and discuss the pros and cons of each designs.
%
\deleted{Our experimental framework and results are available in the open for use by the community.}


%%% ER: not needed for submission, but does not hurt to keep them
\begin{IEEEkeywords}
Broadcast, MANETs, Experimentation
\end{IEEEkeywords}

\end{abstract}


